\chapter{Trabajos relacionados}

En esta sección se presentan trabajos relacionados con el uso de modelos de aprendizaje profundo para la detección de enfermedades en
plantas. La revisión se organiza en dos grupos principales. En primer lugar, se analizan investigaciones que comparan distintas
arquitecturas de redes neuronales convolucionales, con especial atención en modelos ligeros como MobileNet y EfficientNet, debido a que
estas familias de modelos fueron utilizadas como base para esta tesis. En segundo lugar, se revisan trabajos que llevaron los modelos
entrenados a aplicaciones móviles, especialmente en Android, con el fin de identificar qué aspectos de la ejecución en dispositivo han
sido evaluados y cuáles permanecen menos explorados.

Esta organización permite ubicar el aporte de la presente investigación. Por un lado, se retoma la comparación entre arquitecturas
ligeras para clasificación de enfermedades en hojas. Por otro lado, se amplía el análisis hacia la ejecución real en dispositivos
móviles, considerando métricas como latencia, memoria, batería, temperatura, tiempo de inicialización y comportamiento frente a imágenes
capturadas en condiciones menos controladas.

\section{Comparación de arquitecturas ligeras para clasificación de enfermedades}

Gopi y Kondaveeti \cite{agronomy13040961} evaluaron 15 arquitecturas convolucionales preentrenadas mediante aprendizaje por
transferencia para el reconocimiento automático de enfermedades en hojas de arroz. El estudio utilizó un conjunto de 10.080 imágenes
organizadas en 10 clases, correspondientes a nueve enfermedades foliares y una clase de hojas sanas. Los modelos fueron entrenados con
técnicas de aumento de datos y evaluados bajo distintos esquemas de partición del conjunto de datos.

Dentro de las arquitecturas evaluadas, EfficientNetB0 y MobileNetV2 resultan especialmente relevantes para esta investigación, ya que
representan familias de modelos orientadas a mantener una relación favorable entre precisión y eficiencia. EfficientNetB0 obtuvo uno de
los mejores desempeños reportados por el estudio, con una exactitud promedio de 99,56 \%, mientras que MobileNetV2 alcanzó una exactitud
promedio de 99,53 \%. Estos resultados muestran que ambas arquitecturas pueden alcanzar un rendimiento alto en tareas de clasificación
foliar cuando son entrenadas sobre conjuntos de datos controlados.

\begin{table}[H]
    \centering
    \caption{Resultados reportados por Gopi y Kondaveeti para MobileNetV2 y EfficientNetB0.}
    \label{tab:related_gopi_kondaveeti}
    \begin{tabular}{|l|c|c|}
        \hline
        \textbf{Métrica} & \textbf{MobileNetV2} & \textbf{EfficientNetB0} \\
        \hline
        Exactitud & 99,53 \% & 99,56 \% \\
        \hline
        Precisión & 97,78 \% & 97,93 \% \\
        \hline
        Sensibilidad & 97,68 \% & 97,80 \% \\
        \hline
        F1-score & 97,68 \% & 97,82 \% \\
        \hline
        Especificidad & 99,74 \% & 99,75 \% \\
        \hline
    \end{tabular}
\end{table}

Aunque los resultados de ambos modelos fueron cercanos, EfficientNetB0 presentó una ligera ventaja en las métricas principales. Sin
embargo, el estudio se concentra en el desempeño del modelo durante la clasificación y no profundiza en la ejecución directa sobre un
dispositivo móvil. Esta diferencia es importante para la presente tesis, ya que un modelo con alta precisión no necesariamente resulta
el más adecuado cuando debe ejecutarse localmente en un teléfono con restricciones de memoria, procesamiento y energía.

De forma similar, Subramaniam et al. \cite{subramaniam2026comparative} realizaron un análisis comparativo de modelos de aprendizaje
profundo para la detección de enfermedades en cultivos. En este trabajo se evaluaron arquitecturas como EfficientNetB0, MobileNet,
InceptionV3, ResNet101 y una CNN personalizada. EfficientNetB0 obtuvo el mejor desempeño del conjunto, alcanzando una precisión de
validación de 97,10 \% y valores cercanos a 0,97 en precisión, sensibilidad y F1-score. Los autores relacionan este desempeño con el
escalado compuesto de EfficientNet, que ajusta de manera conjunta profundidad, ancho y resolución de entrada.

En contraste, MobileNet presentó el rendimiento más bajo entre los modelos comparados, con una precisión de validación de 94,20 \%.
Este resultado muestra una tensión que también es relevante para esta investigación: las arquitecturas más ligeras pueden ser adecuadas
para escenarios con recursos limitados, pero su menor complejidad puede reducir la capacidad para capturar patrones visuales más
difíciles. Por esta razón, la comparación entre MobileNetV3 Small y EfficientNetV2B0 no solo debe considerar la exactitud, sino también
el costo práctico de ejecutar cada modelo dentro de una aplicación móvil.

\section{Aplicaciones móviles para diagnóstico de enfermedades en plantas}

Ahmed y Reddy \cite{ahmed2021mobile} propusieron un sistema móvil para la detección de enfermedades en hojas de plantas basado en redes
neuronales convolucionales. El sistema fue implementado como una aplicación Android capaz de capturar imágenes y clasificar 38 categorías
asociadas a 14 especies de cultivos, utilizando un conjunto de datos que incluye imágenes de PlantVillage. El trabajo reporta una
precisión de clasificación de 94 \% y un tiempo promedio de predicción de 0,88 segundos, mostrando la viabilidad de integrar modelos de
aprendizaje profundo dentro de una aplicación móvil.

Este trabajo es relevante porque vincula directamente el entrenamiento de modelos con una aplicación Android orientada al diagnóstico de
enfermedades foliares. Sin embargo, el análisis se centra principalmente en la precisión y en el tiempo de predicción, sin presentar una
evaluación más amplia del comportamiento del modelo en el dispositivo. No se reportan métricas como consumo de memoria, batería,
temperatura, estabilidad entre sesiones o diferencias de rendimiento entre dispositivos con distintas capacidades de hardware.

Niaz et al. \cite{Niaz2025Efficient} desarrollaron una aplicación móvil para la detección de enfermedades en trigo, orientada a su uso en
campo por agricultores. A diferencia de los trabajos centrados únicamente en imágenes, este estudio utilizó un conjunto de datos textual
curado a partir de fuentes gubernamentales, expertos y literatura agrícola. Se evaluaron distintos algoritmos de aprendizaje automático,
incluyendo árboles de decisión, bosques aleatorios, máquinas de soporte vectorial y AdaBoost, junto con técnicas de extracción de
características como Count Vectorizer y TF-IDF. El mejor resultado fue obtenido por SVM con Count Vectorizer, alcanzando una precisión de
99 \%.

El aporte principal de este trabajo se encuentra en la construcción de una herramienta móvil orientada al diagnóstico en campo, con un modo
de funcionamiento \textit{offline} para zonas con conectividad limitada. Esta característica se relaciona con el objetivo de la presente
tesis, que también busca una solución capaz de funcionar sin depender de servicios externos. Sin embargo, al igual que en otros trabajos
orientados a aplicaciones móviles, la evaluación se concentra en métricas del modelo y no incluye un análisis detallado de la ejecución
\textit{on-device}, como latencia en distintos dispositivos, memoria utilizada, consumo de batería o temperatura.

\section{Síntesis de los trabajos revisados}

Los trabajos revisados muestran que las arquitecturas ligeras como MobileNet y EfficientNet han sido utilizadas con buenos resultados en
tareas de clasificación de enfermedades foliares. En general, EfficientNet tiende a obtener mejores métricas de clasificación, mientras
que MobileNet resulta atractiva por su diseño más liviano y su orientación a entornos con recursos limitados. Esta relación entre
precisión y eficiencia justifica la selección de MobileNetV3 Small y EfficientNetV2B0 como arquitecturas comparativas en esta
investigación.

También se observa que existen propuestas que integran modelos de aprendizaje automático o aprendizaje profundo en aplicaciones Android
para diagnóstico agrícola. Sin embargo, la mayoría de estos trabajos reporta principalmente métricas de clasificación del modelo, y en
menor medida analiza el comportamiento técnico de la aplicación una vez que el modelo son ejecutados en los dispositivo. 
